\hypertarget{orientation}{}
\hypertarget{orientation-what-you-will-be-able-to-do}{%
\chapter{Orientation: what you will be able to
do}\label{orientation-what-you-will-be-able-to-do}}

By the end, a student should be able to reproduce every analytical task
in the paper, explain why it is needed, distinguish a valid claim from
an overclaim, and design a leakage-free portfolio benchmarking pipeline
from raw monthly data to an auditable report. This orientation chapter
sets the scope, states measurable learning outcomes, previews the data
and the headline results, maps the chapters, lists the prerequisites and
conventions, and warns of the misconceptions the course is built to
prevent. Read it once before starting and return to the roadmap and the
misconception table whenever a later chapter feels disconnected from the
whole.

\hypertarget{what-these-notes-are-about}{%
\section{What these notes are about}\label{what-these-notes-are-about}}

These notes teach a single, self-contained empirical study end to end:
the risk-adjusted benchmarking of the 25 Fama--French size and
book-to-market portfolios. The object of study is deliberately modest ---
25 passively constructed, value-weighted characteristic portfolios and
three factor return series --- because the difficulty lies not in the
data but in \emph{what may honestly be concluded from it}. The same 25
portfolios are the classical test assets from which the size and value
factors were themselves built, so their factor ``alpha'' is a
model-relative residual, not a measure of skill.

The intellectual arc is one long argument. We first \textbf{describe}
the historical cross-section of alpha; then \textbf{quantify} how much of
the apparent ordering could be sampling noise; then test whether the
ordering \textbf{survives} resampling, time and a change of benchmark;
and finally ask whether a ranking formed from the past carries genuine
\textbf{predictive} content for the future. The punchline, developed
carefully across the chapters, is a distinction that is easy to state and
easy to get wrong: a \emph{stable} cross-sectional ordering does carry
out-of-sample information, but \emph{re-estimating} that ordering through
time --- the thing practitioners usually call ``dynamic'' skill --- adds
nothing here. Learning to see, defend and not overstate that distinction
is the point of the course.

Two commitments run through every chapter. The first is
\textbf{uncertainty-first}: a number is never reported without an honest
statement of how much it could move under resampling, model choice or
time. The second is \textbf{reproducibility}: every figure, table and
claim is regenerated by a deterministic pipeline from hashed inputs, so a
reader can audit rather than trust.

\hypertarget{learning-outcomes}{%
\section{Learning outcomes}\label{learning-outcomes}}

On completing these notes a student should be able to:

\begin{enumerate}
\tightlist
\item
  \textbf{Estimate.} Fit a factor regression for each portfolio by
  ordinary least squares, and explain precisely why the intercept is
  called alpha and what it does \emph{not} measure.
\item
  \textbf{Quantify dependence.} Build the full $25\times25$
  heteroskedasticity- and autocorrelation-consistent (HAC) covariance of
  the alpha vector, and explain why 25 separate standard errors are
  insufficient when the portfolios share dates and shocks.
\item
  \textbf{Pool.} Derive and apply multivariate empirical-Bayes shrinkage,
  and articulate the winner's-curse problem it addresses.
\item
  \textbf{Test.} Carry out the joint HAC/Wald test and the Gibbons--Ross--Shanken
  test, control the false-discovery rate across 25 hypotheses, and say
  which test is primary and why.
\item
  \textbf{Resample.} Implement a common-date circular block bootstrap
  that refits every stage, and read a rank interval correctly as a
  resampling distribution rather than a Gaussian confidence interval.
\item
  \textbf{Validate forward.} Construct a strictly leakage-free forward
  evaluation with frozen training estimates, and prove that no future
  information enters a score, rank or loading.
\item
  \textbf{Decompose.} Run the incremental test that separates the
  dynamic content of a rolling ranking from the static content of a fixed
  ordering, and recognise it as a Diebold--Mariano test of equal
  predictive accuracy.
\item
  \textbf{Interpret.} Separate a supported claim from an overclaim, and
  write a conclusion that survives adversarial reading.
\end{enumerate}

The verbs are deliberate: the assessment (Chapter~18) asks a student to
\emph{do} each of these, not merely to recognise them.

\hypertarget{the-dataset-at-a-glance}{%
\section{The dataset at a glance}\label{the-dataset-at-a-glance}}

\begin{longtable}[]{@{}p{0.34\textwidth}p{0.60\textwidth}@{}}
\toprule\noalign{}
Quantity & Value in these notes \\
\midrule\noalign{}
\endhead
\bottomrule\noalign{}
\endlastfoot
Cross-section $N$ & 25 size $\times$ book-to-market portfolios \\
Months $T$ & 1{,}193 consecutive months, July 1926 -- November 2025 \\
Panel size & 29{,}825 portfolio-months, balanced, no calendar gaps \\
Benchmark factors $K$ & 3 for FF3 (Mkt$-$RF, SMB, HML); 5 for the FF5 sensitivity \\
Return units & Decimal monthly excess return, $y_{it}=r_{it}-r_{ft}$ \\
Annualisation & Monthly decimal alpha $\times\,12\times10{,}000$ = annualised basis points \\
Bootstrap & 5{,}000 circular 12-month blocks, seed 2026 \\
Rolling design & 120-month training windows, stepped 12 months; 12-month forward horizon \\
\end{longtable}

The panel is validated by a \emph{fail-closed} loader (Chapter~4): it
raises rather than proceeds if it finds a duplicate portfolio-month key,
a null required field, within-month factor disagreement, an unbalanced
panel or a missing calendar month. Data integrity is treated as part of
the statistical specification, not as pre-processing hygiene.

\hypertarget{headline-results-you-will-be-able-to-defend}{%
\section{Headline results you will be able to
defend}\label{headline-results-you-will-be-able-to-defend}}

The following numbers are the destination. Each is derived, audited and
qualified in a later chapter; they are collected here so that the
machinery always has a visible purpose.

\begin{longtable}[]{@{}p{0.46\textwidth}p{0.48\textwidth}@{}}
\toprule\noalign{}
Finding & Figure \\
\midrule\noalign{}
\endhead
\bottomrule\noalign{}
\endlastfoot
Joint zero-alpha hypothesis rejected (HAC/Wald) & $p = 6.4\times10^{-10}$;
GRS $p = 1.2\times10^{-7}$ (secondary) \\
Best and worst posterior alpha & \texttt{SMALL\ HiBM} $+132$; \texttt{SMALL\ LoBM}
$-185$ bps/year \\
Individually significant alphas after FDR & 5 of 25 survive 5\% Benjamini--Hochberg \\
Point ranking is not a precise hierarchy & bootstrap rank intervals wide;
only 2 of 25 posterior intervals exclude zero \\
Long-horizon mobility & stay-in-quintile probability $\approx 25\%$ over
120 months (near the $20\%$ of pure mobility) \\
Forward quintile spread (top minus bottom) & $243$ bps/year, HAC 95\% CI
$[153,\,333]$, positive in $73\%$ of windows \\
Incremental (rolling minus expanding) skill & $-0.042$, HAC $t=-1.87$:
\emph{no} dynamic edge \\
Benchmark dependence & \texttt{ME5\ BM2} moves rank 9 (FF3) to rank 20 (FF5) \\
Stochastic-frontier support & only 1 of 25 boundary tests survives FDR \\
Data provenance & local vs.\ current official snapshot differ; classified
\textsc{not\_comparable}, not \textsc{pass} \\
\end{longtable}

The single most important line is the incremental result: the positive
forward spread is real, but it is explained by a fixed, persistent
characteristic pattern, not by dynamic timing.

\hypertarget{how-the-chapters-fit-together}{%
\section{How the chapters fit
together}\label{how-the-chapters-fit-together}}

\begin{longtable}[]{@{}p{0.06\textwidth}p{0.42\textwidth}p{0.46\textwidth}@{}}
\toprule\noalign{}
Ch. & Topic & What you should be able to do afterwards \\
\midrule\noalign{}
\endhead
\bottomrule\noalign{}
\endlastfoot
1 & Orientation & State the argument, the data and the conventions \\
2 & Preliminaries: mathematical foundations & Derive every tool used later
--- OLS, HAC, empirical Bayes, the bootstrap, Diebold--Mariano, and more \\
3 & The research question & Frame the estimand and its seven difficulties \\
4 & Data, construction, units, provenance & Build a validated balanced
panel and separate replication from reconciliation \\
5 & Factor regression and the meaning of alpha & Derive OLS alpha and
interpret it as model-relative \\
6 & Full cross-portfolio Newey--West HAC & Assemble and audit the joint
$25\times25$ covariance \\
7 & Joint testing and false discovery & Run HAC/Wald and GRS; control the
FDR across 25 tests \\
8 & Multivariate empirical Bayes & Derive and apply partial pooling \\
9 & Circular block bootstrap & Produce rank and score intervals by full
refit \\
10 & Rolling persistence and transitions & Diagnose the overlap illusion;
read mobility \\
11 & Strictly forward validation & Freeze training estimates and prove no
leakage \\
12 & The incremental test & Separate dynamic forecast from stable ordering \\
13 & Residual diagnostics and model risk & Justify HAC and FF3/FF5
sensitivity \\
14 & Stochastic frontier analysis & Derive it fully and see why it stays
a narrow diagnostic \\
15 & External validation and status & Assign a defensible validation
status \\
16 & Reproducible implementation & Run the deterministic pipeline
end to end \\
17 & Reading the paper's evidence & Map each claim to its supporting layer \\
18 & Capstone & Reproduce the complete study to an assessment standard \\
19 & Exercises with guided solutions & Practise every derivation and audit \\
20 & Terminology glossary & Use each term precisely \\
21 & References and further reading & Trace every method to a source \\
\end{longtable}

\hypertarget{prerequisites}{%
\section{Prerequisites}\label{prerequisites}}

The course is written for a numerate reader --- a BSc mathematics
graduate is the target --- but it is self-contained wherever a finance
concept is introduced. The prerequisites below are the tools assumed
without re-derivation; everything specific to asset pricing is built up
in the chapters.

\hypertarget{mathematics}{%
\subsection{Mathematics}\label{mathematics}}

Linear algebra: vectors and matrices, the transpose, matrix products,
symmetric and positive (semi-)definite matrices, the eigen-decomposition,
matrix inverses and, more importantly, why one solves a linear system
rather than forming an inverse. Probability: expectation, variance and
covariance as operators; the multivariate normal density and its
quadratic form; conditional and marginal normals. Analysis and
optimisation: gradients, first-order conditions, completing the square,
and one-dimensional numerical optimisation. Asymptotics at the level of
the law of large numbers and the central limit theorem, enough to read
``$\hat\theta$ is consistent and asymptotically normal'' without alarm.
Each of these appears in a derivation: the eigen-decomposition floors a
covariance matrix (Chapter~6), completing the square yields the empirical-Bayes
posterior (Chapter~8), and quadratic forms drive the Wald statistic
(Chapter~7); all are also derived in the preliminaries of Chapter~2.

\hypertarget{statistics}{%
\subsection{Statistics}\label{statistics}}

Ordinary least squares and its finite-sample and asymptotic properties;
the meaning of a standard error as sampling variability of an estimator,
which is \emph{not} the volatility of the underlying returns; the logic
of a hypothesis test and the correct reading of a $p$-value as a
tail probability under a null, never as the probability that the null is
true; confidence intervals as repeated-sampling statements; the idea of
resampling (the bootstrap) to approximate a sampling distribution; and
the multiple-comparisons problem, which motivates false-discovery
control. A reader who has only ever run a single regression and read a
single $t$-statistic will find the leap to \emph{joint} inference over 25
correlated estimates the main conceptual hurdle of the course; it is
worth pausing on.

\hypertarget{computing}{%
\subsection{Computing}\label{computing}}

Python 3 with NumPy, pandas, SciPy, statsmodels and matplotlib;
\texttt{pytest} for tests; and a reproducible command-line workflow
(virtual environments, pinned requirements, a single entry-point script,
and deterministic random seeds). No deep software engineering is assumed,
but the course insists on habits that make results auditable: pinned
dependencies, hashed inputs, a recorded run manifest, and tests that fail
loudly. The reference implementation targets CPython 3.13.

\hypertarget{a-five-minute-self-check}{%
\subsection{A five-minute
self-check}\label{a-five-minute-self-check}}

If the following questions are comfortable, the prerequisites are met; if
two or more are not, the cited chapter fills the gap.

\begin{enumerate}
\tightlist
\item
  Why does $(X'X)^{-1}X'y$ minimise $\lVert y-X\theta\rVert^2$? (Chapter~2)
\item
  If two estimators are positively correlated, is the variance of their
  difference larger or smaller than the sum of their variances?
  (Chapter~2)
\item
  A $p$-value is $0.03$. State one thing it does and one thing it does
  \emph{not} mean. (Chapter~2)
\item
  Why does choosing the largest of 25 noisy estimates bias it upward?
  (Chapter~2)
\item
  Why can a correlation between a past ranking and future returns exist
  even with no dynamic forecasting skill? (Chapters~10 and~12)
\end{enumerate}

\hypertarget{mathematical-conventions-used-throughout}{%
\section{Mathematical conventions used
throughout}\label{mathematical-conventions-used-throughout}}

\begin{longtable}[]{@{}p{0.20\textwidth}p{0.34\textwidth}p{0.36\textwidth}@{}}
\toprule\noalign{}
Notation & Meaning & Example \\
\midrule\noalign{}
\endhead
\bottomrule\noalign{}
\endlastfoot
Lower-case italic & Scalar & \emph{r}\textsubscript{it},
\ensuremath{\alpha }\textsubscript{i}, \ensuremath{\sigma } \\
Bold or context-defined lower-case & Vector & f\textsubscript{t} is K \ensuremath{\times }
1; \ensuremath{\alpha } is N \ensuremath{\times } 1 \\
Upper-case & Matrix or sample size & X is T \ensuremath{\times } (K+1); V is N \ensuremath{\times } N \\
Prime, \ensuremath{{}^{\prime}} & Transpose & X\ensuremath{{}^{\prime}}X \\
Hat, \textasciicircum{} & Estimated quantity & \ensuremath{\hat{\alpha }}, \ensuremath{\hat{\beta }}, \ensuremath{\hat{V}} \\
E(\ensuremath{\cdot }), Var(\ensuremath{\cdot }), Cov(\ensuremath{\cdot }) & Expectation, variance, covariance & Var(\ensuremath{\hat{\alpha }}) \\
N, T, K & Portfolios, months, factors & N=25, T=1,193, K=3 for FF3 \\
$y_{it}$ & Excess return of portfolio $i$ in month $t$ & $r_{it}-r_{ft}$ \\
$f_t$ & Factor return vector in month $t$ & Mkt$-$RF, SMB, HML \\
$\varepsilon_{it}$ & Regression residual & $y_{it}-\hat\alpha_i-\hat\beta_i' f_t$ \\
$V$, $V_{\mathrm{HAC}}$ & Covariance of the alpha vector & $N\times N$, HAC \\
$\mu$, $\tau$ & Empirical-Bayes prior mean and dispersion & learned from the
cross-section \\
$L$ & HAC (Bartlett) lag length & $L=12$ months \\
$A_i^{+}$ & Frozen-beta forward alpha & next 12 months, out of sample \\
$Q_p$ & Empirical $p$-quantile & bootstrap intervals \\
\end{longtable}

All vectors are column vectors unless stated otherwise. Matrix inverses
in equations explain the mathematics; implementations should normally
use linear solves or factorisations rather than explicitly computing an
inverse.

\textbf{Units and signs.} Returns and factors are decimal monthly figures
(a return of one percent is $0.01$, not $1$). A monthly decimal alpha $a$
is reported as $a\times 12\times 10{,}000$ annualised basis points; this
is a linear reporting convention, not geometric compounding, so
$a=0.001$ becomes $120$ bps/year. A \emph{positive} alpha is
benchmark-relative outperformance in-sample; a \emph{negative} alpha is a
shortfall. Ranks run from 1 (highest posterior alpha) to $N=25$ (lowest),
and quintile 5 is the highest-scoring group of five portfolios. High
book-to-market is \emph{value}; low book-to-market is \emph{growth} ---
reversing these is the single most common labelling error in the course.

\hypertarget{recommended-reading-routes}{%
\section{Recommended reading routes}\label{recommended-reading-routes}}

\hypertarget{first-pass}{%
\subsection{First pass (concepts)}\label{first-pass}}

Chapters 3--5, 7, 10--12 and 17. Concentrate on the research question, the
meaning of alpha, the joint tests, the persistence and forward chapters,
the incremental decomposition, and how the paper's evidence is read.
Skip the derivations on this pass and aim to state each conclusion and
its main caveat in one sentence.

\hypertarget{implementation-pass}{%
\subsection{Implementation pass (code and
audit)}\label{implementation-pass}}

Chapters 4--9 and 15--18. Reproduce the pipeline stage by stage --- panel
construction, OLS, the joint HAC covariance, the tests, empirical Bayes
and the bootstrap --- then the validation, the full reproducible
implementation and the capstone. Audit every intermediate output against
the numerical checks the chapters provide.

\hypertarget{advanced-pass}{%
\subsection{Advanced pass (theory)}\label{advanced-pass}}

Start with the mathematical foundations in Chapter~2, then the applied
derivations that build on them: the joint HAC derivation (Chapter~6), the
empirical-Bayes posterior and hyperparameter derivations (Chapter~8), the
GRS statistic (Chapter~7), block-bootstrap theory (Chapter~9) and the
complete stochastic-frontier chapter (Chapter~14). This pass fills in
every ``why this estimator'' left implicit on the first two passes.

\hypertarget{the-governing-principle}{%
\section{The governing principle}\label{the-governing-principle}}

\textbf{Estimate, quantify uncertainty, validate forward, then
interpret.} A ranking is not evidence merely because it can be computed.
The research object is not the ordered list alone; it is the ordered
list together with sampling uncertainty, benchmark dependence, temporal
stability, and leakage-free predictive evidence.

The order of the four verbs is not cosmetic. Estimation without
uncertainty invites a false hierarchy; uncertainty without forward
validation cannot distinguish a persistent pattern from an in-sample
artefact; and forward validation without a careful incremental decomposition
cannot tell a genuinely \emph{dynamic} forecast from a
\emph{stable} ordering that any fixed ranking would have captured.
Interpretation comes last precisely because it is the step at which
overclaiming is most tempting and least excusable. A recurring discipline
follows from the principle: every headline number is paired with the
answer to ``how much could this move, and under what?''.

\hypertarget{four-layers-of-evidence}{%
\section{Four layers of evidence}\label{four-layers-of-evidence}}

\begin{longtable}[]{@{}p{0.20\textwidth}p{0.34\textwidth}p{0.36\textwidth}@{}}
\toprule\noalign{}
Layer & Question & Typical tools \\
\midrule\noalign{}
\endhead
\bottomrule\noalign{}
\endlastfoot
Description & What happened in the historical sample? & Raw returns,
factor alpha, point ranks \\
Inference & Could the observed pattern be sampling noise? & HAC
covariance, Wald/GRS, FDR, posterior intervals \\
Stability & Does the conclusion survive resampling, time and model
choice? & Block bootstrap, persistence, transition matrices, FF3/FF5
sensitivity \\
Prediction & Does information available at time t help after time t? &
Frozen training estimates, forward windows, incremental tests \\
\end{longtable}

The layers are cumulative, not alternative. A claim earns credibility
only by surviving each layer in turn: a portfolio with a large historical
alpha (description) whose posterior interval excludes zero (inference),
whose rank is stable under resampling and across FF3/FF5 (stability), and
whose ordering carries information into strictly future data
(prediction). Most of the 25 portfolios fail at the second layer; the
cross-section as a whole passes the joint test but, as the incremental
chapter shows, the predictive layer supports only the weaker of two
possible conclusions.

\hypertarget{common-misconceptions-this-course-prevents}{%
\section{Common misconceptions this course
prevents}\label{common-misconceptions-this-course-prevents}}

\begin{longtable}[]{@{}p{0.46\textwidth}p{0.48\textwidth}@{}}
\toprule\noalign{}
Tempting but wrong & The disciplined statement \\
\midrule\noalign{}
\endhead
\bottomrule\noalign{}
\endlastfoot
``Alpha measures skill.'' & Alpha is the return the \emph{chosen}
benchmark fails to explain; on these passive test assets it is a pricing
residual, not skill. \\
``The top-ranked portfolio is best.'' & Bootstrap rank intervals are wide
and overlapping; the point rank is not a precise hierarchy. \\
``Adjacent rolling windows agree, so performance persists.'' & Adjacent
120-month windows overlap by 90\%; the agreement is mechanical. Only
non-overlapping horizons measure structural persistence. \\
``Past ranking predicts future returns, so the ranking has skill.'' & A
\emph{stable} ordering predicts too; the incremental test shows rolling
re-estimation adds nothing. \\
``GRS rejects, so the result is solid.'' & GRS assumes IID-normal errors,
which the diagnostics reject; the HAC/Wald test is primary here. \\
``Twenty-five $t$-tests at 5\% are fine.'' & Twenty-five tests inflate
false discoveries; the false-discovery rate must be controlled. \\
``The five-factor result should match the three-factor result.'' & Alpha
is benchmark-dependent; a portfolio can move 11 ranks between FF3 and
FF5. \\
``Our numbers match the official data.'' & They do not match the current
official snapshot exactly; the comparison is \textsc{not\_comparable}
until the vintage is verified. \\
``The 243 bps spread is a trading profit.'' & It is a gross,
pre-cost research diagnostic that shorts small-growth; it is not an
implementable return. \\
\end{longtable}

\hypertarget{software-environment-and-reproducibility}{%
\section{Software environment and
reproducibility}\label{software-environment-and-reproducibility}}

Every result in these notes is regenerated by a deterministic pipeline
from two hashed input files. The reviewed environment is CPython 3.13
with pinned versions of NumPy, pandas, SciPy, statsmodels and matplotlib;
the bootstrap uses the fixed seed 2026; and a run manifest records the
input SHA-256 digests, the configuration and the runtime so that a reader
can confirm exactly what produced a figure. The canonical sequence is to
install the pinned requirements, run the test suite, run the analysis
pipeline, export the tables, and build the report. Two properties make
the pipeline auditable rather than merely runnable: the loader is
\emph{fail-closed} (it raises on any integrity violation rather than
silently proceeding), and the local input files are treated as
\emph{immutable} and identified by hash, so a reference download can never
quietly change the analytical sample. Chapters~15 and~16 make these
mechanisms concrete.

\hypertarget{how-to-use-these-notes}{%
\section{How to use these
notes}\label{how-to-use-these-notes}}

Work with a terminal and an editor open beside the text; the notes are
meant to be executed, not only read. When a chapter states a number,
regenerate it and confirm it matches before moving on --- the habit of
auditing is half of what the course teaches. Attempt each exercise in
Chapter~19 before reading its guided solution, and treat the capstone in
Chapter~18 as the real assessment: a complete, leakage-free reproduction
of the study with a written conclusion that neither overstates nor
undersells the evidence. Keep two questions in view throughout: \emph{how
uncertain is this?} and \emph{would this claim survive an adversarial
reader?} If a sentence cannot answer both, it is not yet a result.

